\chapter{Introduction}
\label{ch:introduction}

\section{Background and Motivation}

The proliferation of connected devices in the Internet of Things (IoT) era has
fundamentally changed the expectations placed on wireless networking infrastructure.
Traditional access-point-centric wireless networks, where all communication is mediated
through a fixed access point, struggle to scale gracefully to large numbers of devices
distributed over wide areas. They create single points of failure, require expensive
cabling infrastructure, and cannot self-organise or self-heal in the event of component
failure.

Wireless Mesh Networks (WMNs) address these limitations by distributing the routing
and forwarding functions across all participating nodes. In a WMN, every device acts
simultaneously as a data source, a relay, and --- when connected to the wired backbone
--- a gateway. Packets traverse the network over multiple wireless hops, following paths
computed by a dedicated routing protocol. This multi-hop architecture provides geographic
coverage far beyond the reach of any single radio transmitter, allows the network to
re-route around failed or congested links, and can be deployed rapidly without
pre-existing infrastructure.

The routing protocol is, consequently, the central design choice in any WMN deployment.
It determines how routes are discovered and maintained, how much overhead is generated
in control traffic, how quickly the network responds to topological changes, and how
energy is consumed at each node. The two broad families of routing protocols that have
been most widely studied are:

\begin{itemize}
  \item \textbf{Proactive (table-driven) protocols}, which continuously maintain routing
    tables for all destinations, trading high control overhead for immediately-available routes.
  \item \textbf{Reactive (on-demand) protocols}, which discover routes only when a
    communication session requires one, trading an initial discovery delay for lower
    steady-state control overhead.
\end{itemize}

The Destination-Sequenced Distance-Vector (DSDV) protocol~\cite{perkins1994dsdv} is the
canonical proactive protocol, extending the classic Bellman-Ford algorithm with sequence
numbers to eliminate routing loops. The Ad hoc On-Demand Distance Vector (AODV)
protocol~\cite{perkins2003aodv}, standardised as RFC~3561, is the canonical reactive
protocol and has been widely deployed in both research and production MANET and WMN systems.

Despite extensive theoretical and simulation-based comparisons in the literature, many
studies are conducted under idealised channel models, simplified traffic patterns, or
without rigorous control of the MAC-layer parameters that determine how control and data
traffic compete for the shared wireless medium. This gap motivates the present work: a
carefully controlled simulation study using a high-fidelity IEEE~802.11 MAC and
physical-layer model, a realistic many-to-one traffic pattern representative of IoT data
collection, and a systematic sweep over network density that reveals the crossover
behaviour between the two protocols.

\section{Problem Statement}

Selecting a routing protocol for a WMN deployment involves navigating a fundamental
tension between two competing design objectives:

\begin{enumerate}
  \item \textbf{Route availability.} A node must have a valid route to the destination
    before it can transmit data. In sparse networks, routes may not yet be available
    (reactive protocols) or may have gone stale (proactive protocols with convergence lag).
  \item \textbf{Channel efficiency.} The wireless medium is shared. Every byte of routing
    control traffic consumes airtime that could otherwise carry application data. In dense
    networks, proactive protocols that broadcast routing tables periodically can saturate
    the medium, leading to MAC-layer congestion and cascading packet loss.
\end{enumerate}

The precise point at which the dominant bottleneck shifts from route availability to
channel efficiency --- the \emph{crossover density} --- is a function of the specific
protocol parameters, radio technology, traffic load, and network size. This crossover is
poorly characterised in the existing literature for the specific combination of AODV,
DSDV, IEEE~802.11g, and many-to-one IoT traffic patterns.

Furthermore, most existing studies report only application-layer performance metrics
(throughput, delay, packet delivery ratio) without tracing the physical mechanism
responsible for observed performance degradation at the MAC layer. This makes it
difficult for practitioners to reason about whether a performance problem is caused by
sub-optimal routing logic or by low-level channel saturation --- an important distinction
when diagnosing and tuning deployed systems.

\section{Research Objectives}

This thesis pursues four specific research objectives:

\begin{enumerate}
  \item \textbf{Quantify the crossover density.} Determine the value of $N$ at which
    AODV begins to outperform DSDV across the five performance metrics studied:
    throughput, packet count, end-to-end delay, energy consumption, and packet drop rate.
  \item \textbf{Explain the crossover mechanism.} Trace the causal chain from routing
    protocol behaviour, through MAC-layer contention, to observed application-layer
    performance, providing quantitative evidence at each level.
  \item \textbf{Provide MAC-layer evidence.} Extract and analyse contention-window
    escalation statistics from the OMNeT++/INET simulation to demonstrate the
    physical-layer mechanism responsible for DSDV's performance collapse at high densities.
  \item \textbf{Derive deployment guidelines.} Translate the quantitative findings into
    actionable engineering guidance for practitioners selecting between AODV and DSDV
    for WMN deployments in IoT and other application domains.
\end{enumerate}

\section{Thesis Contributions}

The principal contributions of this thesis are:

\begin{itemize}
  \item A fully reproducible OMNeT++~6.3 / INET~4.5 simulation setup for comparing
    AODV and DSDV in a WMN context, with all configuration files provided in
    Appendix~\ref{app:config}.
  \item Empirical demonstration of the AODV/DSDV performance crossover at approximately
    $N = 20$ nodes, with quantitative data for throughput, delay, energy, and drop rate
    across all five tested densities.
  \item MAC-layer evidence linking DSDV's performance degradation at $N = 50$ to
    contention-window escalation and the 7-retry abort limit defined by the
    IEEE~802.11 standard.
  \item A comprehensive energy analysis showing AODV's 8\% energy advantage at high
    densities and explaining the reversal at $N = 10$.
  \item Protocol-selection guidelines for WMN practitioners, structured as a decision
    matrix covering network size, latency requirements, reliability requirements, and
    energy constraints.
\end{itemize}

\section{Thesis Structure}

The remainder of this thesis is organised as follows. Chapter~\ref{ch:wmn} provides
background on WMN architecture, characteristics, and applications. Chapter~\ref{ch:routing}
reviews the theoretical foundations of proactive and reactive routing, describes the AODV
and DSDV protocols in detail, and analyses the IEEE~802.11 DCF MAC layer.
Chapter~\ref{ch:methodology} describes the simulation methodology, including the network
model, parameter choices, and metric collection procedures. Chapter~\ref{ch:results}
presents and analyses the simulation results, with particular attention to the trade-offs
between the two protocols. Chapter~\ref{ch:conclusion} summarises the findings, provides
deployment guidelines, and outlines directions for future work.
